% Generated mechanically from frozen input inputs/p12/ja/stacks-perfect.tex.
% Do not edit this generated file; edit the bound locale source instead.

% OK, start here.
%

\setcounter{chapter}{103}
\chapter{代数スタックの導来圏}


\phantomsection
\label{stacks-perfect-section-phantom}





\section{はじめに}
\label{stacks-perfect-section-introduction}

\noindent
本章では，代数スタックに付随する導来圏について述べる．
とくに準連接層の導来圏を扱う．すなわち，スキームに関する結果
（Derived Categories of Schemes, Section
\ref{perfect-section-introduction} を参照）および代数空間に関する結果
（Derived Categories of Spaces, Section
\ref{spaces-perfect-section-introduction} を参照）に対応する類似結果を証明する．
本章の結果が \cite{LM-B} の結果と異なる主な理由は，
一貫して「大サイト」を用いることにある．本章を読む前に，
ここで用いる用語を導入した「代数スタック上の層」および
「代数スタックのコホモロジー」の各章に目を通されたい．



\section{規約・記法と言葉の濫用}
\label{stacks-perfect-section-conventions}

\noindent
Properties of Stacks, Section
\ref{stacks-properties-section-conventions} で導入した規約および
言葉の濫用を引き続き用いる．記法については
Cohomology of Stacks, Section
\ref{stacks-cohomology-section-notation} の説明に従う．











\section{lisse-\'etale サイトと flat-fppf サイト}
\label{stacks-perfect-section-lisse-etale}

\noindent
本節は，導来圏についての Cohomology of Stacks, Section
\ref{stacks-cohomology-section-lisse-etale} の類似である．

\begin{lemma}
\label{stacks-perfect-lemma-shriek-derived}
$\mathcal{X}$ を代数スタックとする．記法は Cohomology of Stacks,
Lemmas \ref{stacks-cohomology-lemma-lisse-etale} および
\ref{stacks-cohomology-lemma-lisse-etale-modules} と同じとする．
\begin{enumerate}
\item 関手
$g_! : \textit{Ab}(\mathcal{X}_{lisse,\etale}) \to
\textit{Ab}(\mathcal{X}_\etale)$
は左導来関手
$$
Lg_! :
D(\mathcal{X}_{lisse,\etale})
\longrightarrow
D(\mathcal{X}_\etale)
$$
をもち，これは $g^{-1}$ の左随伴であり，$g^{-1}Lg_! = \text{id}$ を満たす．
\item 関手 $g_! : 
\textit{Mod}(\mathcal{X}_{lisse,\etale},
\mathcal{O}_{\mathcal{X}_{lisse,\etale}}) \to
\textit{Mod}(\mathcal{X}_\etale, \mathcal{O}_{\mathcal{X}})$
は左導来関手
$$
Lg_! :
D(\mathcal{O}_{\mathcal{X}_{lisse,\etale}})
\longrightarrow
D(\mathcal{X}_\etale, \mathcal{O}_\mathcal{X})
$$
をもち，これは $g^*$ の左随伴であり，$g^*Lg_! = \text{id}$ を満たす．
\item 関手 $g_! : \textit{Ab}(\mathcal{X}_{flat,fppf}) \to
\textit{Ab}(\mathcal{X}_{fppf})$
は左導来関手
$$
Lg_! :
D(\mathcal{X}_{flat, fppf})
\longrightarrow
D(\mathcal{X}_{fppf})
$$
をもち，これは $g^{-1}$ の左随伴であり，$g^{-1}Lg_! = \text{id}$ を満たす．
\item 関手 $g_! :
\textit{Mod}(\mathcal{X}_{flat,fppf},
\mathcal{O}_{\mathcal{X}_{flat,fppf}}) \to
\textit{Mod}(\mathcal{X}_{fppf}, \mathcal{O}_{\mathcal{X}})$
は左導来関手
$$
Lg_! :
D(\mathcal{O}_{\mathcal{X}_{flat, fppf}})
\longrightarrow
D(\mathcal{O}_\mathcal{X})
$$
をもち，これは $g^*$ の左随伴であり，$g^*Lg_! = \text{id}$ を満たす．
\end{enumerate}
注意：加群に対する $Lg_!$ がアーベル層に対する $Lg_!$ と
一致するかどうかは先験的には明らかでない．
Cohomology on Sites, Remark
\ref{sites-cohomology-remark-when-derived-shriek-equal} を参照されたい．
\end{lemma}

\begin{proof}
関手 $Lg_!$ の存在と $g^*$ の左随伴であることは
Cohomology on Sites, Lemma
\ref{sites-cohomology-lemma-existence-derived-lower-shriek} である．
（アーベル層の場合には，定数層 $\mathbf{Z}$ を構造層として用いる．）
さらに，複体 $\mathcal{H}^\bullet$ 上での値は，
適切な左分解
$\mathcal{K}^\bullet \to \mathcal{H}^\bullet$
を取り，$\mathcal{K}^\bullet$ に関手 $g_!$ を適用することで計算される．
Cohomology of Stacks, Lemmas
\ref{stacks-cohomology-lemma-lisse-etale-modules} および
\ref{stacks-cohomology-lemma-lisse-etale} により
$g^{-1}g_!\mathcal{K}^\bullet = \mathcal{K}^\bullet$ であるから，
いずれの場合にも最後の主張が成り立つ．
\end{proof}

\begin{lemma}
\label{stacks-perfect-lemma-lisse-etale-functorial-derived}
Cohomology of Stacks, Lemma
\ref{stacks-cohomology-lemma-lisse-etale-functorial}
と同じ仮定と記法のもとで，
$$
g^{-1} \circ Rf_* = Rf'_* \circ (g')^{-1}
\quad\text{and}\quad
L(g')_! \circ (f')^{-1} = f^{-1} \circ Lg_!
$$
が非有界導来圏上で成り立つ
（加群の場合にもアーベル層の場合にも成り立つ）．
\end{lemma}

\begin{proof}
$\tau = \etale$（それぞれ\ $\tau = fppf$）とする．
$\mathcal{F}$ を $\mathcal{X}_\tau$ 上のアーベル層とする．
Cohomology of Stacks, Lemma
\ref{stacks-cohomology-lemma-lisse-etale-functorial-cohomology}
により，標準的な（基底変換）写像
$$
g^{-1}Rf_*\mathcal{F} \longrightarrow Rf'_*(g')^{-1}\mathcal{F}
$$
は同型である．残りの証明は形式的である．
アーベル群の層として計算したコホモロジーと加群層として計算した
コホモロジーとは一致するので，$\mathcal{F}$ が
$\mathcal{X}_\tau$ 上の加群層である場合にも
$g^{-1} Rf_*\mathcal{F} = Rf'_* (g')^{-1}\mathcal{F}$ を得る．

\medskip\noindent
次に，$\mathcal{Y}_{lisse,\etale}$
（それぞれ\ $\mathcal{Y}_{flat,fppf}$）上の
$\mathcal{G}$（加群層またはアーベル群の層）に対し，標準写像
$$
L(g')_!(f')^{-1}\mathcal{G} \to f^{-1}Lg_!\mathcal{G}
$$
が同型であることを示す．このためには，
$\mathcal{X}_\tau$ 上の任意の単射的層 $\mathcal{I}$ に対して，
誘導される写像
$$
\Hom(L(g')_!(f')^{-1}\mathcal{G}, \mathcal{I}[n])
\leftarrow
\Hom(f^{-1}Lg_!\mathcal{G}, \mathcal{I}[n])
$$
がすべての $n \in \mathbf{Z}$ に対して同型であることを示せば十分である
（Hom は適切な導来圏で取る）．
$f^{-1}$ と $Rf_*$ の随伴性，$Lg_!$ と $g^{-1}$ の随伴性，
およびそれらの「プライム付き」版により，これは上で証明した同型
$g^{-1} Rf_*\mathcal{I} \to Rf'_* (g')^{-1}\mathcal{I}$ から従う．

\medskip\noindent
$\mathcal{Y}_{lisse,\etale}$
（それぞれ\ $\mathcal{Y}_{fppf}$）上の
有界複体 $\mathcal{G}^\bullet$（加群またはアーベル群）については，
標準写像
\begin{equation}
\label{stacks-perfect-equation-to-show}
L(g')_!(f')^{-1}\mathcal{G}^\bullet \to f^{-1}Lg_!\mathcal{G}^\bullet
\end{equation}
は同型である．実際，これは切断を用いる通常の議論と，
$L(g')_!(f')^{-1}$ および $f^{-1}Lg_!$ が三角圏の完全関手であることから，
層の場合に帰着される．

\medskip\noindent
$\mathcal{G}^\bullet$ が
$\mathcal{Y}_{lisse,\etale}$
（それぞれ\ $\mathcal{Y}_{fppf}$）上の
上に有界な複体（加群またはアーベル群）であるとする．
標準写像 (\ref{stacks-perfect-equation-to-show}) は，
stupid 切断 $\sigma_{\geq -n}$（Homology, Section
\ref{homology-section-truncations} を参照）を用いて，
$\mathcal{G}^\bullet$ を有界複体の余極限
$\mathcal{G}^\bullet = \colim \mathcal{G}_n^\bullet$
として書けるため同型である．
これにより識別三角形
$$
\bigoplus\nolimits_{n \geq 1} \mathcal{G}_n^\bullet \to
\bigoplus\nolimits_{n \geq 1} \mathcal{G}_n^\bullet \to
\mathcal{G}^\bullet \to \ldots
$$
を得る．また，関手 $L(g')_!$, $(f')^{-1}$, $f^{-1}$, $Lg_!$ は
いずれも（複体の）直和と可換である．

\medskip\noindent
$\mathcal{G}^\bullet$ が
$\mathcal{Y}_{lisse,\etale}$
（それぞれ\ $\mathcal{Y}_{fppf}$）上の任意の複体
（加群またはアーベル群）である場合には，標準切断
$\tau_{\leq n}$（Homology, Section
\ref{homology-section-truncations} を参照）を用いて
$\mathcal{G}^\bullet$ を上に有界な複体の余極限として書き，
直前の段落の議論を繰り返す．

\medskip\noindent
最後に，$f^{-1}$ と $Rf_*$ の随伴性，$Lg_!$ と $g^{-1}$ の随伴性，
およびそれらの「プライム付き」版により，補題の第一の等式が
完全な一般性のもとで第二の等式から従うことが分かる．
\end{proof}

\begin{lemma}
\label{stacks-perfect-lemma-higher-shriek-quasi-coherent}
$\mathcal{X}$ を代数スタックとする．記法は Cohomology of Stacks,
Lemma \ref{stacks-cohomology-lemma-lisse-etale} と同じものとする．
\begin{enumerate}
\item $\mathcal{H}$ を，$\mathcal{X}$ の lisse-\'etale サイト上の準連接
$\mathcal{O}_{\mathcal{X}_{lisse,\etale}}$-加群とする．
すべての $p \in \mathbf{Z}$ に対し，層 $H^p(Lg_!\mathcal{H})$ は
$\mathcal{X}$ 上で平坦基底変換性をもつ局所的準連接加群である．
\item $\mathcal{H}$ を，$\mathcal{X}$ の flat-fppf サイト上の準連接
$\mathcal{O}_{\mathcal{X}_{flat,fppf}}$-加群とする．
すべての $p \in \mathbf{Z}$ に対し，層 $H^p(Lg_!\mathcal{H})$ は
$\mathcal{X}$ 上で平坦基底変換性をもつ局所的準連接加群である．
\end{enumerate}
\end{lemma}

\begin{proof}
スキーム $U$ と全射的な滑らかな射 $x : U \to \mathcal{X}$ を取る．
Modules on Sites, Definition \ref{sites-modules-definition-site-local}
により，各引き戻し $f_i^{-1}\mathcal{H}$ が大域表示をもつような
\'etale（それぞれ\ fppf）被覆
$\{U_i \to U\}_{i \in I}$ が存在する（Modules on Sites, Definition
\ref{sites-modules-definition-global} を参照）．
ここで $f_i : U_i \to \mathcal{X}$ は合成
$U_i \to U \to \mathcal{X}$ であり，代数スタックの射である．
（引き戻しが $\mathcal{X}/f_i$ への制限「そのもの」であることを
思い出されたい．Sheaves on Stacks, Definition
\ref{stacks-sheaves-definition-pullback} およびその直後の議論を参照．）
被覆を細分することにより，各 $U_i$ がアフィンスキームであると
仮定してよい．各 $f_i$ は滑らか（それぞれ\ 平坦）なので，
Lemma \ref{stacks-perfect-lemma-lisse-etale-functorial-derived} により
$f_i^{-1}Lg_!\mathcal{H} = Lg_{i, !}(f'_i)^{-1}\mathcal{H}$ である．
Cohomology of Stacks, Lemma
\ref{stacks-cohomology-lemma-check-lqc-fbc-on-covering}
を用いると，補題の主張は，$\mathcal{H}$ が大域表示をもち，かつある
アフィンスキーム $X = \Spec(A)$ に対して
$\mathcal{X} = (\Sch/X)_{fppf}$ である場合に帰着される．

\medskip\noindent
この表示を
$$
\bigoplus\nolimits_{j \in J} \mathcal{O} \longrightarrow
\bigoplus\nolimits_{i \in I} \mathcal{O} \longrightarrow
\mathcal{H} \longrightarrow 0
$$
と書く．ここで $\mathcal{O} = \mathcal{O}_{\mathcal{X}_{lisse,\etale}}$
（それぞれ\ $\mathcal{O} = \mathcal{O}_{\mathcal{X}_{flat,fppf}}$）である．
サイト $\mathcal{X}_{lisse,\etale}$
（それぞれ\ $\mathcal{X}_{flat,fppf}$）は終対象，すなわち
準コンパクトな $X/X$ をもつことに注意する
（Cohomology on Sites, Section
\ref{sites-cohomology-section-limits} を参照）．したがって Sites, Lemma
\ref{sites-lemma-directed-colimits-sections} により
$$
\Gamma(\bigoplus\nolimits_{i \in I} \mathcal{O}) =
\bigoplus\nolimits_{i \in I} A
$$
である．ゆえに，この表示の写像は，ある $A$-加群 $M$ の同様な表示
$$
\bigoplus\nolimits_{j \in J} A \longrightarrow
\bigoplus\nolimits_{i \in I} A \longrightarrow
M \longrightarrow 0
$$
に対応する．さらに，$\mathcal{H}$ は $M$ に付随する準連接層 $M^a$ を
lisse-\'etale（それぞれ\ flat-fppf）サイトへ制限したものに等しい．
自由 $A$-加群による分解
$$
\ldots \to F_2 \to F_1 \to F_0 \to M \to 0
$$
を選ぶ．$\mathcal{X}_{lisse,\etale}$
（それぞれ\ $\mathcal{X}_{flat,fppf}$）の各対象 $U/X$ に対して
構造射 $U \to X$ は平坦なので，複体
$$
\ldots \mathcal{O} \otimes_A F_2 \to \mathcal{O} \otimes_A F_1 \to
\mathcal{O} \otimes_A F_0 \to \mathcal{H} \to 0
$$
は自由 $\mathcal{O}$-加群による $\mathcal{H}$ の分解である．
したがって構成により，$Lg_!\mathcal{H}$ の値は
$$
\ldots \to
\mathcal{O}_\mathcal{X} \otimes_A F_2 \to
\mathcal{O}_\mathcal{X} \otimes_A F_1 \to
\mathcal{O}_\mathcal{X} \otimes_A F_0 \to 0 \to \ldots
$$
である．これは $\mathcal{X}_\etale$
（それぞれ\ $\mathcal{X}_{fppf}$）上の準連接加群の複体なので，
Cohomology of Stacks, Proposition
\ref{stacks-cohomology-proposition-loc-qcoh-flat-base-change}
から $H^p(Lg_!\mathcal{H})$ は準連接である．
\end{proof}

\section{コホモロジーと lisse-\'etale および flat-fppf サイト}
\label{stacks-perfect-section-compare-unbounded-lisse-etale}

\noindent
代数スタック $\mathcal{X}$ 上の層のコホモロジーは flat-fppf サイト上で
計算できることをすでに見た．この節では，$\mathcal{X}$ の direct 圏の
（有界とは限らない）対象についても同じことが成り立つと証明する．
% P12-ERR-0041: frozen source says "direct category"; preserve literally pending canon disposition.

\begin{lemma}
\label{stacks-perfect-lemma-higher-shriek-Z}
$\mathcal{X}$ を代数スタックとする．Lemma
\ref{stacks-perfect-lemma-shriek-derived} の (1) の $Lg_!$ についても，Lemma
\ref{stacks-perfect-lemma-shriek-derived} の (3) の $Lg_!$ についても
$Lg_!\mathbf{Z} = \mathbf{Z}$ が成り立つ．
\end{lemma}

\begin{proof}
flat-fppf サイトと fppf サイトとの比較について証明する．
lisse-\'etale サイトの場合も全く同じである．
$i \not = 0$ に対して $H^i(Lg_!\mathbf{Z})$ が $0$ であり，かつ標準写像
$H^0(Lg_!\mathbf{Z}) \to \mathbf{Z}$ が同型であることを示せばよい．
$\mathcal{U}$ がスキームであり，さらに $f$ が局所有限表示であるような，
全射的平坦射 $f : \mathcal{U} \to \mathcal{X}$ を取る．
（例えば表示 $U \to \mathcal{X}$ を取り，$\mathcal{U}$ を $U$ に対応する
代数スタックとすればよい．）Sheaves on Stacks, Lemmas
\ref{stacks-sheaves-lemma-check-exactness-covering} および
\ref{stacks-sheaves-lemma-surjective-flat-locally-finite-presentation}
により，$i \not = 0$ に対して引き戻し
$f^{-1}H^i(Lg_!\mathbf{Z})$ が $0$ であり，かつ引き戻し
$H^0(Lg_!\mathbf{Z}) \to f^{-1}\mathbf{Z}$ が同型であることを
示せば十分である．Lemma
\ref{stacks-perfect-lemma-lisse-etale-functorial-derived} により
$f^{-1}Lg_!\mathbf{Z} = L(g')_!\mathbf{Z}$ である．ここで
$g' : \Sh(\mathcal{U}_{flat, fppf}) \to \Sh(\mathcal{U}_{fppf})$
は $\mathcal{U}$ に対する対応する比較射である．
これで次の段落で扱う場合に帰着された．

\medskip\noindent
あるスキーム $X$ に対して
$\mathcal{X} = (\Sch/X)_{fppf}$ であると仮定する．
この場合，圏 $\mathcal{X}_{flat, fppf}$ は終対象 $e$，すなわち $X/X$ をもち，
さらに関手 $u : \mathcal{X}_{flat, fppf} \to \mathcal{X}_{fppf}$ は
$e$ を終対象へ写す．終対象が存在するとき，$\mathbf{Z}$ は終対象上の
自由アーベル層である．したがって Cohomology on Sites, Lemma
\ref{sites-cohomology-lemma-existence-derived-lower-shriek}
における $Lg_!$ の構成そのものから
$Lg_!\mathbf{Z} = \mathbf{Z}$ を得る．
\end{proof}

\begin{lemma}
\label{stacks-perfect-lemma-lisse-etale-cohomology}
$\mathcal{X}$ を代数スタックとする．記法は Lemma
\ref{stacks-perfect-lemma-shriek-derived} と同じものとする．
\begin{enumerate}
\item $D(\mathcal{X}_\etale)$ の $K$ に対して，
\begin{enumerate}
\item $R\Gamma(\mathcal{X}_\etale, K) =
R\Gamma(\mathcal{X}_{lisse,\etale}, g^{-1}K)$ であり，かつ
\item $\mathcal{X}_{lisse,\etale}$ の任意の対象 $x$ に対して
$R\Gamma(x, K) =
R\Gamma(\mathcal{X}_{lisse,\etale}/x, g^{-1}K)$ である．
\end{enumerate}
\item $D(\mathcal{X}_{fppf})$ の $K$ に対して，
\begin{enumerate}
\item $R\Gamma(\mathcal{X}_{fppf}, K) =
R\Gamma(\mathcal{X}_{flat,fppf}, g^{-1}K)$ であり，かつ
\item $\mathcal{X}_{flat,fppf}$ の任意の対象 $x$ に対して
$H^p(x, K) =
R\Gamma(\mathcal{X}_{flat,fppf}/x, g^{-1}K)$ である．
% P12-ERR-0043: frozen source equates H^p with a derived object; preserve literally pending canon disposition.
\end{enumerate}
\end{enumerate}
いずれの場合にも加群について同じことが成り立つ．実際，
$g^{-1} = g^*$ であり，Cohomology on Sites, Lemma
\ref{sites-cohomology-lemma-modules-abelian-unbounded}
によりコホモロジーの計算に相違はない．
\end{lemma}

\begin{proof}
flat-fppf サイトと fppf サイトとの比較について証明する．
lisse-\'etale サイトの場合も全く同じである．Lemma
\ref{stacks-perfect-lemma-higher-shriek-Z} により $Lg_!\mathbf{Z} = \mathbf{Z}$ である．
したがって
\begin{align*}
R\Gamma(\mathcal{X}_{fppf}, K)
& =
R\Hom(\mathbf{Z}, K) \\
& =
R\Hom(Lg_!\mathbf{Z}, K) \\
& =
R\Hom(\mathbf{Z}, g^{-1}K) \\
& =
R\Gamma(\mathcal{X}_{lisse,\etale}, g^{-1}K)
\end{align*}
を得る．
% P12-ERR-0044: frozen flat-fppf comparison ends at lisse-etale; preserve literally pending canon disposition.
これで (1)(a) が証明された．
(1)(b) は (1)(a) から従う．実際，$x$ がスキーム $U$ 上にあるなら，
サイト $\mathcal{X}_\etale/x$ は $(\Sch/U)_\etale$ と同値であり，
$\mathcal{X}_{lisse,\etale}$ は $U_{lisse, \etale}$ と同値である．
% P12-ERR-0045/P12-ERR-0046: frozen proof retains (1) numbering and etale sites inside the stated flat-fppf case.
\end{proof}

\section{準連接加群の導来圏}
\label{stacks-perfect-section-derived}

\noindent
$\mathcal{X}$ を代数スタックとする．包含関手
$\QCoh(\mathcal{O}_\mathcal{X}) \to
\textit{Mod}(\mathcal{O}_\mathcal{X})$ は完全でないので，
$D_\QCoh(\mathcal{O}_\mathcal{X})$ を，準連接なコホモロジー層をもつ
複体からなる $D(\mathcal{O}_\mathcal{X})$ の充満部分圏として
定義することはできない．そこで Cohomology of Stacks, Remark
\ref{stacks-cohomology-remark-bousfield-colocalization}
との類推により，準連接加群の導来圏を商として定義する．

\medskip\noindent
$\textit{LQCoh}^{fbc}(\mathcal{O}_\mathcal{X}) \subset
\textit{Mod}(\mathcal{O}_\mathcal{X})$ は，平坦基底変換性をもつ
局所的準連接 $\mathcal{O}_\mathcal{X}$-加群の充満部分圏を表すことを
思い出そう．Cohomology of Stacks, Section
\ref{stacks-cohomology-section-loc-qcoh-flat-base-change} を参照されたい．
以下では
$$
D_{\textit{LQCoh}^{fbc}}(\mathcal{O}_\mathcal{X}) =
D_{\textit{LQCoh}^{fbc}(\mathcal{O}_\mathcal{X})}(\mathcal{O}_\mathcal{X})
$$
と略記する．Derived Categories, Lemma
\ref{derived-lemma-cohomology-in-serre-subcategory} および
Cohomology of Stacks, Proposition
\ref{stacks-cohomology-proposition-loc-qcoh-flat-base-change} の (2) から，
$D_{\textit{LQCoh}^{fbc}}(\mathcal{O}_\mathcal{X})$ は
$D(\mathcal{O}_\mathcal{X})$ の厳密充満な飽和三角部分圏であると分かる．

\medskip\noindent
$\textit{Parasitic}(\mathcal{O}_\mathcal{X}) \subset
\textit{Mod}(\mathcal{O}_\mathcal{X})$ を，寄生的な
$\mathcal{O}_\mathcal{X}$-加群の充満部分圏とする．
Cohomology of Stacks, Section
\ref{stacks-cohomology-section-parasitic} を参照されたい．以下では
$$
D_{\textit{Parasitic}}(\mathcal{O}_\mathcal{X}) =
D_{\textit{Parasitic}(\mathcal{O}_\mathcal{X})}(\mathcal{O}_\mathcal{X})
$$
と略記する．$\textit{Parasitic}(\mathcal{O}_\mathcal{X})$ は
$\textit{Mod}(\mathcal{O}_\mathcal{X})$ の Serre 部分圏なので，
先と同様に，これは $D(\mathcal{O}_\mathcal{X})$ の厳密充満な
飽和三角部分圏である．Cohomology of Stacks, Lemma
\ref{stacks-cohomology-lemma-parasitic} を参照されたい．

\medskip\noindent
$\textit{Mod}(\mathcal{O}_\mathcal{X})$ の弱 Serre 部分圏の共通部分
$\textit{Parasitic}(\mathcal{O}_\mathcal{X}) \cap
\textit{LQCoh}^{fbc}(\mathcal{O}_\mathcal{X})$
もまた弱 Serre 部分圏である．同様に
\begin{align*}
D_{\textit{Parasitic} \cap \textit{LQCoh}^{fbc}}(\mathcal{O}_\mathcal{X})
& =
D_{\textit{Parasitic}(\mathcal{O}_\mathcal{X}) \cap
\textit{LQCoh}^{fbc}(\mathcal{O}_\mathcal{X})}(\mathcal{O}_\mathcal{X}) \\
& =
D_{\textit{Parasitic}}(\mathcal{O}_\mathcal{X})
\cap
D_{\textit{LQCoh}^{fbc}}(\mathcal{O}_\mathcal{X})
\end{align*}
と略記する．先と同様に，これは $D(\mathcal{O}_\mathcal{X})$ の
厳密充満な飽和三角部分圏である．したがってなおさら，
$D_{\textit{LQCoh}^{fbc}}(\mathcal{O}_\mathcal{X})$ の
厳密充満な飽和三角部分圏でもある．

\begin{definition}
\label{stacks-perfect-definition-derived}
$\mathcal{X}$ を代数スタックとする．上の記法のもとで，
{\it 準連接なコホモロジー層をもつ $\mathcal{O}_\mathcal{X}$-加群の導来圏}
を Verdier 商\footnote{この定義は文献における定義とは異なる．
\cite[6.3]{olsson_sheaves} を参照されたい．しかし Lemma
\ref{stacks-perfect-lemma-derived-quasi-coherent} により，この定義はその定義と一致する．}
$$
D_\QCoh(\mathcal{O}_\mathcal{X}) =
D_{\textit{LQCoh}^{fbc}}(\mathcal{O}_\mathcal{X})/
D_{\textit{Parasitic} \cap \textit{LQCoh}^{fbc}}(\mathcal{O}_\mathcal{X})
$$
として定義する．
\end{definition}

\noindent
Verdier 商は Derived Categories, Section
\ref{derived-section-quotients} で定義されている．
$D_{\textit{LQCoh}^{fbc}}(\mathcal{O}_\mathcal{X})$ の射
$a : E \to E'$ が $D_\QCoh(\mathcal{O}_\mathcal{X})$ で同型になるための
必要十分条件は，錐 $C(a)$ が寄生的なコホモロジー層をもつことである．
Derived Categories, Lemma \ref{derived-lemma-operations} を参照されたい．

\medskip\noindent
関手
$$
D_{\textit{LQCoh}^{fbc}}(\mathcal{O}_\mathcal{X})
\xrightarrow{H^i}
\textit{LQCoh}^{fbc}(\mathcal{O}_\mathcal{X})
\xrightarrow{Q}
\QCoh(\mathcal{O}_\mathcal{X})
$$
を考える．$Q$ は部分圏
$\textit{Parasitic}(\mathcal{O}_\mathcal{X}) \cap
\textit{LQCoh}^{fbc}(\mathcal{O}_\mathcal{X})$ を零化することに注意する．
Cohomology of Stacks, Lemma
\ref{stacks-cohomology-lemma-adjoint-kernel-parasitic} を参照されたい．
Derived Categories, Lemma
\ref{derived-lemma-universal-property-quotient} により，コホモロジー的関手
\begin{equation}
\label{stacks-perfect-equation-Hi-quasi-coherent}
H^i :
D_\QCoh(\mathcal{O}_\mathcal{X})
\longrightarrow
\QCoh(\mathcal{O}_\mathcal{X})
\end{equation}
を得る．さらに，$E \in D_\QCoh(\mathcal{O}_\mathcal{X})$ が零であるための
必要十分条件は，すべての $i \in \mathbf{Z}$ に対して
$H^i(E) = 0$ であることにも注意する．実際 Cohomology of Stacks,
Lemma \ref{stacks-cohomology-lemma-adjoint-kernel-parasitic} により，
$Q$ の核はちょうど
$\textit{Parasitic}(\mathcal{O}_\mathcal{X}) \cap
\textit{LQCoh}^{fbc}(\mathcal{O}_\mathcal{X})$ に等しい．

\medskip\noindent
圏 $\textit{Parasitic}(\mathcal{O}_\mathcal{X}) \cap
\textit{LQCoh}^{fbc}(\mathcal{O}_\mathcal{X})$ および
$\textit{LQCoh}^{fbc}(\mathcal{O}_\mathcal{X})$ は，\'etale 位相における
加群のアーベル圏
$\textit{Mod}(\mathcal{X}_\etale, \mathcal{O}_\mathcal{X})$
の弱 Serre 部分圏でもあることに注意する．Cohomology of Stacks,
Proposition \ref{stacks-cohomology-proposition-loc-qcoh-flat-base-change}
および Lemma \ref{stacks-cohomology-lemma-parasitic} を参照されたい．
したがって次の補題の主張は意味をもつ．

\begin{lemma}
\label{stacks-perfect-lemma-compare-etale-fppf-QCoh}
$\mathcal{X}$ を代数スタックとする．
$\mathcal{P}_\mathcal{X} = \textit{Parasitic}(\mathcal{O}_\mathcal{X}) \cap
\textit{LQCoh}^{fbc}(\mathcal{O}_\mathcal{X})$ と略記する．比較射
$\epsilon : \mathcal{X}_{fppf} \to \mathcal{X}_\etale$ は可換図式
$$
\xymatrix{
D_{\textit{Parasitic} \cap \textit{LQCoh}^{fbc}}(\mathcal{O}_\mathcal{X})
\ar[r] &
D_{\textit{LQCoh}^{fbc}}(\mathcal{O}_\mathcal{X}) \ar[r] &
D(\mathcal{O}_\mathcal{X}) \\
D_{\mathcal{P}_\mathcal{X}}(\mathcal{X}_\etale, \mathcal{O}_\mathcal{X})
\ar[r] \ar[u]^{\epsilon^*} &
D_{\textit{LQCoh}^{fbc}(\mathcal{O}_\mathcal{X})}(
\mathcal{X}_\etale, \mathcal{O}_\mathcal{X})
\ar[r] \ar[u]^{\epsilon^*} &
D(\mathcal{X}_\etale, \mathcal{O}_\mathcal{X})
\ar[u]^{\epsilon^*}
}
$$
を誘導する．さらに左の二つの垂直矢印は三角圏の同値である．したがって
$$
D_{\textit{LQCoh}^{fbc}(\mathcal{O}_\mathcal{X})}
(\mathcal{X}_\etale, \mathcal{O}_\mathcal{X})
/
D_{\mathcal{P}_\mathcal{X}}(\mathcal{X}_\etale, \mathcal{O}_\mathcal{X})
\longrightarrow
D_\QCoh(\mathcal{O}_\mathcal{X})
$$
という同値も得られる．
\end{lemma}

\begin{proof}
$\epsilon^*$ は完全なので，補題の主張にある図式が得られることは明らかである．
Cohomology on Sites, Lemma
\ref{sites-cohomology-lemma-compare-topologies-derived-adequate-modules}
を次の状況に適用して，中央の垂直矢印が同値であることを示す：
$\mathcal{C} = \mathcal{X}$，
$\tau = fppf$，
$\tau' = \etale$，
$\mathcal{O} = \mathcal{O}_\mathcal{X}$，
$\mathcal{A} = \textit{LQCoh}^{fbc}(\mathcal{O}_\mathcal{X})$，そして
$\mathcal{B}$ はアフィンスキーム上にある $\mathcal{X}$ の対象の集合とする．
この補題を適用できることを確認するには，条件 (1), (2), (3), (4) を
調べなければならない．条件 (1) と (2) は上の議論から明らかである
（明示的には Cohomology of Stacks, Proposition
\ref{stacks-cohomology-proposition-loc-qcoh-flat-base-change} から従う）．
すべてのスキームはアフィン開集合による Zariski 開被覆をもつので，
条件 (3) が成り立つ．条件 (4) は Descent, Lemma
\ref{descent-lemma-quasi-coherent-and-flat-base-change} から従う．

\medskip\noindent
圏の同値
$\epsilon^* :
D_{\textit{LQCoh}^{fbc}(\mathcal{O}_\mathcal{X})}(
\mathcal{X}_\etale, \mathcal{O}_\mathcal{X})
\to
D_{\textit{LQCoh}^{fbc}}(\mathcal{O}_\mathcal{X})$
が，寄生的なコホモロジー層をもつ複体の部分圏の同値を誘導することの
確認は省略する．
\end{proof}

\noindent
$\mathcal{X}$ を代数スタックとする．Cohomology of Stacks, Lemma
\ref{stacks-cohomology-lemma-quasi-coherent-weak-serre} により，

準連接加群の圏
$\QCoh(\mathcal{O}_{\mathcal{X}_{lisse,\etale}})$ は
$\textit{Mod}(\mathcal{O}_{\mathcal{X}_{lisse,\etale}})$ の
弱 Serre 部分圏をなす．

同様に，準連接加群の圏
$\QCoh(\mathcal{O}_{\mathcal{X}_{flat,fppf}})$ は
$\textit{Mod}(\mathcal{O}_{\mathcal{X}_{flat,fppf}})$ の
弱 Serre 部分圏をなす．したがって
$$
D_\QCoh(\mathcal{O}_{\mathcal{X}_{lisse,\etale}}) =
D_{\QCoh(\mathcal{O}_{\mathcal{X}_{lisse,\etale}})}(
\mathcal{O}_{\mathcal{X}_{lisse,\etale}})
\subset
D(\mathcal{O}_{\mathcal{X}_{lisse,\etale}})
$$
を考えることができ，同様に
$$
D_\QCoh(\mathcal{O}_{\mathcal{X}_{flat,fppf}}) =
D_{\QCoh(\mathcal{O}_{\mathcal{X}_{flat,fppf}})}(
\mathcal{O}_{\mathcal{X}_{flat,fppf}})
\subset
D(\mathcal{O}_{\mathcal{X}_{flat,fppf}})
$$
を考えることができる．先と同様に，これらは厳密充満な飽和三角部分圏である．
$D_\QCoh(\mathcal{O}_\mathcal{X})$ は，実はこれらのいずれとも同値である．

\begin{lemma}
\label{stacks-perfect-lemma-derived-quasi-coherent}
$\mathcal{X}$ を代数スタックとする．
$\mathcal{P}_\mathcal{X} = \textit{Parasitic}(\mathcal{O}_\mathcal{X}) \cap
\textit{LQCoh}^{fbc}(\mathcal{O}_\mathcal{X})$ と置く．
\begin{enumerate}
\item
$\mathcal{F}^\bullet$ を
$D_{\textit{LQCoh}^{fbc}(\mathcal{O}_\mathcal{X})}
(\mathcal{X}_\etale, \mathcal{O}_\mathcal{X})$ の対象とする．
Cohomology of Stacks, Lemma
\ref{stacks-cohomology-lemma-lisse-etale} における lisse-\'etale サイトの
$g$ に対して，次が成り立つ．
\begin{enumerate}
\item $g^*\mathcal{F}^\bullet$ は
$D_\QCoh(\mathcal{O}_{\mathcal{X}_{lisse,\etale}})$ に属する．
\item $g^*\mathcal{F}^\bullet = 0$ であるための必要十分条件は，
$\mathcal{F}^\bullet$ が
$D_{\mathcal{P}_\mathcal{X}}(\mathcal{X}_\etale, \mathcal{O}_\mathcal{X})$
に属することである．
\item $\mathcal{H}^\bullet$ が
$D_\QCoh(\mathcal{O}_{\mathcal{X}_{lisse,\etale}})$ に属するなら，
$Lg_!\mathcal{H}^\bullet$ は
$D_{\textit{LQCoh}^{fbc}(\mathcal{O}_\mathcal{X})}(
\mathcal{X}_\etale, \mathcal{O}_\mathcal{X})$ に属する．
\item 関手 $g^*$ と $Lg_!$ は互いに逆な関手
$$
\xymatrix{
D_\QCoh(\mathcal{O}_\mathcal{X}) \ar@<1ex>[r]^-{g^*} &
D_\QCoh(\mathcal{O}_{\mathcal{X}_{lisse,\etale}})
\ar@<1ex>[l]^-{Lg_!}
}
$$
を定める．
\end{enumerate}
\item
$\mathcal{F}^\bullet$ を
$D_{\textit{LQCoh}^{fbc}}(\mathcal{O}_\mathcal{X})$ の対象とする．
Cohomology of Stacks, Lemma
\ref{stacks-cohomology-lemma-lisse-etale} における flat-fppf サイトの
$g$ に対して，次が成り立つ．
\begin{enumerate}
\item $g^*\mathcal{F}^\bullet$ は
$D_\QCoh(\mathcal{O}_{\mathcal{X}_{flat, fppf}})$ に属する．
\item $g^*\mathcal{F}^\bullet = 0$ であるための必要十分条件は，
$\mathcal{F}^\bullet$ が
$D_{\mathcal{P}_\mathcal{X}}(\mathcal{O}_\mathcal{X})$ に属することである．
\item $\mathcal{H}^\bullet$ が
$D_\QCoh(\mathcal{O}_{\mathcal{X}_{flat,fppf}})$ に属するなら，
$Lg_!\mathcal{H}^\bullet$ は
$D_{\textit{LQCoh}^{fbc}}(\mathcal{O}_\mathcal{X})$ に属する．
\item 関手 $g^*$ と $Lg_!$ は互いに逆な関手
$$
\xymatrix{
D_\QCoh(\mathcal{O}_\mathcal{X}) \ar@<1ex>[r]^-{g^*} &
D_\QCoh(\mathcal{O}_{\mathcal{X}_{flat,fppf}}) \ar@<1ex>[l]^-{Lg_!}
}
$$
を定める．
\end{enumerate}
\end{enumerate}
\end{lemma}

\begin{proof}
関手 $g^* = g^{-1}$ は完全である．したがって (1)(a), (2)(a),
(1)(b), (2)(b) は Cohomology of Stacks, Lemmas
\ref{stacks-cohomology-lemma-quasi-coherent} および
\ref{stacks-cohomology-lemma-parasitic-in-terms-flat-fppf} から従う．

\medskip\noindent
(1)(c) と (2)(c) を証明する．Lemma \ref{stacks-perfect-lemma-shriek-derived}
における $Lg_!$ の構成を用いる．この構成は Cohomology on Sites, Lemma
\ref{sites-cohomology-lemma-existence-derived-lower-shriek} を介する．

さらに Derived Categories, Proposition
\ref{derived-proposition-left-derived-exists} を用いる．これらの構成によれば，
$D(\mathcal{O}_{\mathcal{X}_{lisse,\etale}})$ の任意の対象
$\mathcal{H}^\bullet$ 上の $Lg_!$ は
$$
Lg_!\mathcal{H}^\bullet = \colim g_!\mathcal{K}_n^\bullet =
g_! \colim \mathcal{K}_n^\bullet
$$
（項ごとの余極限）として計算される．ここで擬同型
$\colim \mathcal{K}_n^\bullet \to \mathcal{H}^\bullet$ は擬同型
$\mathcal{K}_n^\bullet \to \tau_{\leq n} \mathcal{H}^\bullet$ を誘導する．
包含関手
$$
\textit{LQCoh}^{fbc}(\mathcal{O}_\mathcal{X}) \subset
\textit{Mod}(\mathcal{X}_\etale, \mathcal{O}_\mathcal{X})
\quad\text{and}\quad
\textit{LQCoh}^{fbc}(\mathcal{O}_\mathcal{X}) \subset
\textit{Mod}(\mathcal{O}_\mathcal{X})
$$
はフィルター余極限と両立する．したがって (c) は
$D_\QCoh(\mathcal{O}_{\mathcal{X}_{lisse,\etale}})$ に属する上に有界な複体，
および $D_\QCoh(\mathcal{O}_{\mathcal{X}_{flat,fppf}})$ に属する上に有界な複体
$\mathcal{H}^\bullet$ について証明すれば十分である．
この場合，$H^n(Lg_!\mathcal{H}^\bullet)$ が
$\textit{LQCoh}^{fbc}(\mathcal{O}_\mathcal{X})$ に属することを示すため，
$i > m$ に対して $\mathcal{H}^i = 0$ となる整数 $m$ に関する帰納法を
用いることができる．$m < n$ なら
$H^n(Lg_!\mathcal{H}^\bullet) = 0$ なので主張は成り立つ．一般の場合，
識別三角形
$$
\tau_{\leq m - 1}\mathcal{H}^\bullet \to \mathcal{H}^\bullet \to
H^m(\mathcal{H}^\bullet)[-m] \to \ldots
$$
（Derived Categories, Remark
\ref{derived-remark-truncation-distinguished-triangle}）を考え，関手
$Lg_!$ を適用する．$\textit{LQCoh}^{fbc}(\mathcal{O}_\mathcal{X})$ は
加群圏の弱 Serre 部分圏なので，三項のうち二項について (c) を
証明すれば十分である．$Lg_!\tau_{\leq m - 1}\mathcal{H}^\bullet$
については帰納法により，$Lg_!H^m(\mathcal{H}^\bullet)[-m]$ については
Lemma \ref{stacks-perfect-lemma-higher-shriek-quasi-coherent} により主張が成り立つ．
したがって (c) が成り立つ．

\medskip\noindent
(2)(d) を証明しよう．(2)(a) と (2)(b) により，関手
$g^{-1} = g^*$ は関手
$$
c :
D_\QCoh(\mathcal{O}_\mathcal{X})
\longrightarrow
D_\QCoh(\mathcal{O}_{\mathcal{X}_{flat, fppf}})
$$
を誘導する．Derived Categories, Lemma
\ref{derived-lemma-universal-property-quotient} を参照されたい．
したがって三角圏の図式
$$
\xymatrix{
D_{\textit{LQCoh}^{fbc}}(\mathcal{O}_\mathcal{X})
\ar[rd]^{g^{-1}} \ar[rr]_q & &
D_\QCoh(\mathcal{O}_\mathcal{X}) \ar[ld]^c \\
& D_\QCoh(\mathcal{O}_{\mathcal{X}_{flat, fppf}})
\ar@<1ex>[lu]^{Lg_!}
}
$$
を得る．ここで $q$ は商関手であり，内側の三角形は可換で，
$g^{-1}Lg_! = \text{id}$ である．

$D_{\textit{LQCoh}^{fbc}}(\mathcal{O}_\mathcal{X})$ の任意の対象 $E$ に対し，
写像 $a : Lg_!g^{-1}E \to E$ は
$D(\mathcal{O}_{\mathcal{X}_{flat, fppf}})$ における擬同型へ写る．
% P12-ERR-0047: frozen source says "For any object of E of"; Japanese supplies ordinary grammar.
したがって $a$ の錐は $g^{-1}$ により零へ写り，(2)(b) により
$q(a)$ は同型である．ゆえに $q \circ Lg_!$ は $c$ の擬逆である．

\medskip\noindent
lisse-\'etale サイトの場合には，上と全く同じ議論により
$$
D_{\textit{LQCoh}^{fbc}(\mathcal{O}_\mathcal{X})}(
\mathcal{X}_\etale, \mathcal{O}_\mathcal{X})
/
D_{\mathcal{P}_\mathcal{X}}(
\mathcal{X}_\etale, \mathcal{O}_\mathcal{X})
$$
が $D_\QCoh(\mathcal{O}_{\mathcal{X}_{lisse,\etale}})$ と同値であると
分かる．Lemma \ref{stacks-perfect-lemma-compare-etale-fppf-QCoh} の最後の同値を
適用すれば証明が完了する．
\end{proof}

\noindent
次の補題は，商関手
$D_{\textit{LQCoh}^{fbc}}(\mathcal{O}_\mathcal{X}) \to
D_\QCoh(\mathcal{O}_\mathcal{X})$ が左随伴をもつことを述べる．
Remark \ref{stacks-perfect-remark-QCoh-admissible} を参照されたい．

\begin{lemma}
\label{stacks-perfect-lemma-bousfield-colocalization}
$\mathcal{X}$ を代数スタックとする．
$E$ を $D_{\textit{LQCoh}^{fbc}}(\mathcal{O}_\mathcal{X})$ の対象とする．
$D_{\textit{LQCoh}^{fbc}}(\mathcal{O}_\mathcal{X})$ には標準的な識別三角形
$$
E' \to E \to P \to E'[1]
$$
が存在し，$P$ は
$D_{\textit{Parasitic} \cap \textit{LQCoh}^{fbc}}
(\mathcal{O}_\mathcal{X})$ に属し，かつ
$$
\Hom_{D(\mathcal{O}_\mathcal{X})}(E', P') = 0
$$
がすべての
$D_{\textit{Parasitic} \cap \textit{LQCoh}^{fbc}}(\mathcal{O}_\mathcal{X})$
の $P'$ に対して成り立つ．
\end{lemma}

\begin{proof}
Cohomology of Stacks, Section
\ref{stacks-cohomology-section-lisse-etale} で考察した環付きトポスの射
$g : \Sh(\mathcal{X}_{flat, fppf}) \longrightarrow \Sh(\mathcal{X}_{fppf})$
を考える．$E' = Lg_!g^*E$ と置き，随伴写像 $E' \to E$ の錐を $P$ とする．
Lemma \ref{stacks-perfect-lemma-shriek-derived} の (4) を参照されたい．Lemma
\ref{stacks-perfect-lemma-derived-quasi-coherent} の (2)(a) および (2)(c) により，
$E'$ は $D_{\textit{LQCoh}^{fbc}}(\mathcal{O}_\mathcal{X})$ に属する．
したがって $P$ も $D_{\textit{LQCoh}^{fbc}}(\mathcal{O}_\mathcal{X})$ に属する．
Lemma \ref{stacks-perfect-lemma-shriek-derived} の (4) により
$g^*Lg_! = \text{id}$ なので，写像 $g^*E' \to g^*E$ は同型である．
よって $g^*P = 0$ であり，Lemma
\ref{stacks-perfect-lemma-derived-quasi-coherent} の (2)(b) により，$P$ は
$D_{\textit{Parasitic} \cap \textit{LQCoh}^{fbc}}(\mathcal{O}_\mathcal{X})$
の対象である．最後に，
$D_{\textit{Parasitic} \cap \textit{LQCoh}^{fbc}}(\mathcal{O}_\mathcal{X})$
の $P'$ に対して
$$
\Hom(E', P') = \Hom(Lg_!g^*E, P') = \Hom(g^*E, g^*P') = 0
$$
である．実際，Lemma \ref{stacks-perfect-lemma-derived-quasi-coherent} の (2)(b) により
$g^*P' = 0$ である．識別三角形
$E' \to E \to P \to E'[1]$ は標準的である．より正確には，$E$ 上の
恒等射を誘導する三角形の同型を除いて一意である．
% P12-ERR-0048: Japanese supplies the missing relative connection in the frozen phrase "isomorphism of triangles induces".
Derived Categories, Section \ref{derived-section-admissible} の議論を
参照されたい．
\end{proof}

\begin{remark}
\label{stacks-perfect-remark-QCoh-admissible}
Lemma \ref{stacks-perfect-lemma-bousfield-colocalization} の結果によれば，
$$
D_{\textit{Parasitic} \cap \textit{LQCoh}^{fbc}}(\mathcal{O}_\mathcal{X})
\subset
D_{\textit{LQCoh}^{fbc}}(\mathcal{O}_\mathcal{X})
$$
は左許容部分圏である．Derived Categories, Section
\ref{derived-section-admissible} を参照されたい．

特に，
$\mathcal{A} \subset D_{\textit{LQCoh}^{fbc}}(\mathcal{O}_\mathcal{X})$
をその左直交補とすると，Derived Categories, Proposition
\ref{derived-proposition-summarize-admissible} により，$\mathcal{A}$ は
$D_{\textit{LQCoh}^{fbc}}(\mathcal{O}_\mathcal{X})$ の右許容部分圏であり，
合成
$$
\mathcal{A} \longrightarrow
D_{\textit{LQCoh}^{fbc}}(\mathcal{O}_\mathcal{X}) \longrightarrow
D_\QCoh(\mathcal{O}_\mathcal{X})
$$
は同値である．すなわち $D_\QCoh(\mathcal{O}_\mathcal{X})$ を
$D_{\textit{LQCoh}^{fbc}}(\mathcal{O}_\mathcal{X})$ の，さらには
$D(\mathcal{X}_{fppf}, \mathcal{O}_\mathcal{X})$ の
厳密充満な飽和三角部分圏とみなすことができる．
\end{remark}

\section{準連接加群の導来順像}
\label{stacks-perfect-section-derived-pushforward}

\noindent
上の内容の最初の応用として導来順像を構成する．Examples, Section
\ref{examples-section-derived-push-quasi-coherent} には，代数スタックの
準コンパクトかつ準分離な射 $f : \mathcal{X} \to \mathcal{Y}$ であって，
順像関手 $Rf_*$ が関手
$D_\QCoh(\mathcal{O}_\mathcal{X}) \to
D_\QCoh(\mathcal{O}_\mathcal{Y})$ を誘導しない例がある．
したがって下に有界な複体へ制限する必要がある．

\begin{proposition}
\label{stacks-perfect-proposition-derived-direct-image-quasi-coherent}
$f : \mathcal{X} \to \mathcal{Y}$ を代数スタックの準コンパクトかつ
準分離な射とする．関手 $Rf_*$ は可換図式
$$
\xymatrix{
D^{+}_{\textit{Parasitic} \cap \textit{LQCoh}^{fbc}}(\mathcal{O}_\mathcal{X})
\ar[r] \ar[d]^{Rf_*} &
D^{+}_{\textit{LQCoh}^{fbc}}(\mathcal{O}_\mathcal{X})
\ar[r] \ar[d]^{Rf_*} &
D(\mathcal{O}_\mathcal{X})
\ar[d]^{Rf_*} \\
D^{+}_{\textit{Parasitic} \cap \textit{LQCoh}^{fbc}}(\mathcal{O}_\mathcal{Y})
\ar[r] &
D^{+}_{\textit{LQCoh}^{fbc}}(\mathcal{O}_\mathcal{Y}) \ar[r] &
D(\mathcal{O}_\mathcal{Y})
}
$$
を誘導し，したがって商圏上の関手
$$
Rf_{\QCoh, *} :
D^{+}_\QCoh(\mathcal{O}_\mathcal{X})
\longrightarrow
D^{+}_\QCoh(\mathcal{O}_\mathcal{Y})
$$
を誘導する．さらに Cohomology of Stacks, Proposition
\ref{stacks-cohomology-proposition-direct-image-quasi-coherent}
の関手 $R^if_\QCoh$ は，(\ref{stacks-perfect-equation-Hi-quasi-coherent}) の $H^i$ に対する
$H^i \circ Rf_{\QCoh, *}$ と等しい．
% P12-ERR-0049: Japanese supplies singular agreement for the frozen "functor ... are equal" wording.
\end{proposition}

\begin{proof}
$D^{+}_{\textit{LQCoh}^{fbc}}(\mathcal{O}_\mathcal{X})$ の $E$ に対して，
$Rf_*E$ が
$D^{+}_{\textit{LQCoh}^{fbc}}(\mathcal{O}_\mathcal{Y})$ の対象であることを
示さなければならない．これは Cohomology of Stacks, Proposition
\ref{stacks-cohomology-proposition-loc-qcoh-flat-base-change} とスペクトル系列
$R^if_*H^j(E) \Rightarrow R^{i + j}f_*E$ から従う．
寄生加群の場合も Cohomology of Stacks, Lemma
\ref{stacks-cohomology-lemma-pushforward-parasitic} を用いて同様に示される．
最後の主張は (\ref{stacks-perfect-equation-Hi-quasi-coherent}) における $H^i$ の定義から
明らかである．
\end{proof}

\section{準連接加群の導来引き戻し}
\label{stacks-perfect-section-derived-pullback}

\noindent
準連接なコホモロジー層をもつ複体の導来引き戻しは一般に存在する．

\begin{proposition}
\label{stacks-perfect-proposition-derived-pullback-quasi-coherent}
$f : \mathcal{X} \to \mathcal{Y}$ を代数スタックの射とする．
完全関手 $f^*$ は可換図式
$$
\xymatrix{
D_{\textit{LQCoh}^{fbc}}(\mathcal{O}_\mathcal{X}) \ar[r] &
D(\mathcal{O}_\mathcal{X}) \\
D_{\textit{LQCoh}^{fbc}}(\mathcal{O}_\mathcal{Y})
\ar[r] \ar[u]^{f^*} &
D(\mathcal{O}_\mathcal{Y}) \ar[u]^{f^*}
}
$$
を誘導する．合成
$$
D_{\textit{LQCoh}^{fbc}}(\mathcal{O}_\mathcal{Y})
\xrightarrow{f^*}
D_{\textit{LQCoh}^{fbc}}(\mathcal{O}_\mathcal{X})
\xrightarrow{q_\mathcal{X}}
D_\QCoh(\mathcal{O}_\mathcal{X})
$$
は局所化
$D_{\textit{LQCoh}^{fbc}}(\mathcal{O}_\mathcal{Y}) \to
D_\QCoh(\mathcal{O}_\mathcal{Y})$ に関して左導来可能であり，
$Lf^*_\QCoh$ をその左導来関手として定義してよい：
$$
Lf_\QCoh^* :
D_\QCoh(\mathcal{O}_\mathcal{Y})
\longrightarrow
D_\QCoh(\mathcal{O}_\mathcal{X})
$$
と定義してよい（Derived Categories, Definitions
\ref{derived-definition-right-derived-functor-defined} および
\ref{derived-definition-everywhere-defined} を参照）．
$f$ が準コンパクトかつ準分離なら，$Lf^*_\QCoh$ と $Rf_{\QCoh, *}$ は
次の随伴性を満たす：
$$
\Hom_{D_\QCoh(\mathcal{O}_\mathcal{X})}(Lf^*_\QCoh A, B)
=
\Hom_{D_\QCoh(\mathcal{O}_\mathcal{Y})}(A, Rf_{\QCoh, *}B)
$$
ただし $A \in D_\QCoh(\mathcal{O}_\mathcal{Y})$，
$B \in D^{+}_\QCoh(\mathcal{O}_\mathcal{X})$ とする．
\end{proposition}

\begin{proof}
最初の主張を証明するには，
$D_{\textit{LQCoh}^{fbc}}(\mathcal{O}_\mathcal{Y})$ の $E$ に対して
$f^*E$ が $D_{\textit{LQCoh}^{fbc}}(\mathcal{O}_\mathcal{X})$ の対象である
ことを示さなければならない．$f^* = f^{-1}$ は完全であり，
Cohomology of Stacks, Proposition
\ref{stacks-cohomology-proposition-loc-qcoh-flat-base-change} により
$f^*$ は $\textit{LQCoh}^{fbc}(\mathcal{O}_\mathcal{Y})$ を
$\textit{LQCoh}^{fbc}(\mathcal{O}_\mathcal{X})$ へ写すので，直ちに従う．

\medskip\noindent
$\mathcal{D} = D_{\textit{LQCoh}^{fbc}}(\mathcal{O}_\mathcal{Y})$ と置く．
$S$ を，錐が
$D_{\textit{Parasitic} \cap \textit{LQCoh}^{fbc}}(\mathcal{O}_\mathcal{Y})$
の対象となる $\mathcal{D}$ の射の集まりとする．
$\mathcal{D}' = D_\QCoh(\mathcal{O}_\mathcal{X})$ と置き，
$F = q_\mathcal{X} \circ f^* : \mathcal{D} \to \mathcal{D}'$ と置く．
すると $\mathcal{D}, S, \mathcal{D}', F$ は Derived Categories,
Situation \ref{derived-situation-derived-functor} および Definition
\ref{derived-definition-right-derived-functor-defined} の状況にある．
$\mathcal{D}$ の任意の対象 $E$ に対して $LF(E)$ が定義されることを
証明しよう．すなわち Lemma \ref{stacks-perfect-lemma-bousfield-colocalization} で構成した
三角形
$$
E' \to E \to P \to E'[1]
$$
を考える．$s : E' \to E$ は $S$ の元であることに注意する．
$E'$ が $LF$ を計算すると主張する．実際，$s' : E'' \to E$ を $S$ の
別の元とする．すなわちこれは，$P'$ が
$D_{\textit{Parasitic} \cap \textit{LQCoh}^{fbc}}(\mathcal{O}_\mathcal{Y})$
に属する三角形
$E'' \to E \to P' \to E''[1]$ に入るとする．Lemma
\ref{stacks-perfect-lemma-bousfield-colocalization}（およびその証明）により，
$E' \to E$ は $E'' \to E$ を経由する．したがって $E' \to E$ は
系 $S/E$ で共終である．ゆえに $E'$ が $LF$ を計算することは明らかである．

\medskip\noindent
最後の主張を示すため，$B = q_\mathcal{X}(H)$ および
$A = q_\mathcal{Y}(E)$ と書く．上のように $E' \to E$ を選ぶ．一方で
$Rf_{\QCoh, *}(B) = q_\mathcal{Y}(Rf_*H)$ を用い，他方で
$Lf^*_\QCoh(A) = q_\mathcal{X}(f^*E')$ を用いる．
\begin{align*}
\Hom_{D_\QCoh(\mathcal{O}_\mathcal{X})}(Lf^*_\QCoh A, B)
& = 
\Hom_{D_\QCoh(\mathcal{O}_\mathcal{X})}(q_\mathcal{X}(f^*E'),
q_\mathcal{X}(H)) \\
& = 
\colim_{H \to H'} \Hom_{D(\mathcal{O}_\mathcal{X})}(f^*E', H') \\
& = \colim_{H \to H'} \Hom_{D(\mathcal{O}_\mathcal{Y})}(E', Rf_*H') \\
& = \Hom_{D(\mathcal{O}_\mathcal{Y})}(E', Rf_*H) \\
& =
\Hom_{D_\QCoh(\mathcal{O}_\mathcal{Y})}(A, Rf_{\QCoh, *}B)
\end{align*}
ここで余極限は，錐 $P(s)$ が
$D^+_{\textit{Parasitic} \cap \textit{LQCoh}^{fbc}}(\mathcal{O}_\mathcal{X})$
の対象となる
$D^+_{\textit{LQCoh}^{fbc}}(\mathcal{O}_\mathcal{X})$ の射
$s : H \to H'$ 全体にわたって取る．第一の等式は上ですでに見た．
第二の等式は Verdier 商の構成から成り立つ．第三の等式は
Cohomology on Sites, Lemma \ref{sites-cohomology-lemma-adjoint} から成り立つ．
Proposition \ref{stacks-perfect-proposition-derived-direct-image-quasi-coherent} により
$Rf_*P(s)$ は
$D^+_{\textit{Parasitic} \cap \textit{LQCoh}^{fbc}}(\mathcal{O}_\mathcal{Y})$
の対象なので，
$\Hom_{D(\mathcal{O}_\mathcal{Y})}(E', Rf_*P(s)) = 0$ である．
したがって第四の等式が成り立つ．最後の等式は $E'$ の構成から成り立つ．
\end{proof}

\section{導来圏における準連接対象}
\label{stacks-perfect-section-QC}

\noindent
この節は Sheaves on Stacks, Section
\ref{stacks-sheaves-section-QC} の続きである．
$\mathcal{X}$ を代数スタックとする．同節では三角圏
$$
\mathit{QC}(\mathcal{X}) = \mathit{QC}(\mathcal{X}_{affine}, \mathcal{O})
$$
を定義し，$\mathcal{X}$ が代数空間 $X$ によって表現可能なら
$\mathit{QC}(\mathcal{X})$ は $D_\QCoh(\mathcal{O}_X)$ と同値であることを
証明した．任意の代数スタックについて同じことを証明するのに十分な理論が，
ここまででちょうど整ったことになる．

\begin{lemma}
\label{stacks-perfect-lemma-cohomology-parasitic}
$\mathcal{X}$ を代数スタックとする．$K$ を，コホモロジー層が寄生的である
$D(\mathcal{X}_{fppf})$ の対象とする．スキーム $U$ 上にあり，
$U \to \mathcal{X}$ が平坦であるような $\mathcal{X}$ のすべての対象 $x$ に
対して $R\Gamma(x, K) = 0$ である．
\end{lemma}

\begin{proof}
Section \ref{stacks-perfect-section-lisse-etale} で論じたトポスの射を
$g : \Sh(\mathcal{X}_{flat, fppf}) \to \Sh(\mathcal{X}_{fppf})$ と書く．
% P12-ERR-0537: Japanese supplies the missing "by/as" relation in both frozen "Denote g ... the morphism" clauses.
$x$ を，スキーム $U$ 上にあり $U \to \mathcal{X}$ が平坦であるような
$\mathcal{X}$ の対象とする．すなわち $x$ は
$\mathcal{X}_{flat, fppf}$ の対象である．
Lemma \ref{stacks-perfect-lemma-lisse-etale-cohomology} の (2)(b) により
$R\Gamma(x, K) = R\Gamma(\mathcal{X}_{flat, fppf}/x, g^{-1}K)$ である．
一方，仮定は $D(\mathcal{X}_{flat, fppf})$ の対象 $g^{-1}K$ の
コホモロジー層が零であることを意味する．Cohomology of Stacks,
Definition \ref{stacks-cohomology-definition-parasitic} を参照されたい．
したがって $g^{-1}K = 0$ であり，証明が完了する．
\end{proof}

\begin{lemma}
\label{stacks-perfect-lemma-cohomology-parasitic-complex}
$\mathcal{X}$ を代数スタックとする．$K$ を $D(\mathcal{X}_{fppf})$ の対象とし，
アフィンスキーム $U$ 上にあり $U \to \mathcal{X}$ が平坦であるような
$\mathcal{X}$ のすべての対象 $x$ に対して $R\Gamma(x, K) = 0$ と仮定する．
このとき，すべての $i$ に対して $H^i(\mathcal{X}, K) = 0$ である．
\end{lemma}

\begin{proof}
Section \ref{stacks-perfect-section-lisse-etale} で論じたトポスの射を
$g : \Sh(\mathcal{X}_{flat, fppf}) \to \Sh(\mathcal{X}_{fppf})$ と書く．
Lemma \ref{stacks-perfect-lemma-lisse-etale-cohomology} の (2)(b) により，仮定は
$g^{-1}K$ が，アフィンスキーム上にある
$\mathcal{X}_{flat, fppf}$ のすべての対象上で消滅コホモロジーをもつことを
意味する．$\mathcal{X}_{flat, fppf}$ のすべての対象 $x$ はそのような対象に
よる被覆をもつので，$g^{-1}K$ のコホモロジー層は消滅する，すなわち
$g^{-1}K = 0$ である．するともちろん
$R\Gamma(\mathcal{X}_{flat, fppf}, g^{-1}K) = 0$ であり，Lemma
\ref{stacks-perfect-lemma-lisse-etale-cohomology} の (2)(a) により求める結論が従う．
\end{proof}

\begin{lemma}
\label{stacks-perfect-lemma-higher-shriek-QC}
$\mathcal{X}$ を代数スタックとする．$K$ を
$D_\QCoh(\mathcal{O}_{\mathcal{X}_{flat, fppf}})$ の対象とする．
このとき $Lg_!K$ は次の性質を満たす：
$\mathcal{X}_{affine}$ の任意の射 $x \to x'$ に対して写像
$$
R\Gamma(x', Lg_!K) \otimes_{\mathcal{O}(x')}^\mathbf{L} \mathcal{O}(x)
\longrightarrow
R\Gamma(x, Lg_!K)
$$
は擬同型である．
\end{lemma}

\begin{proof}
Lemma \ref{stacks-perfect-lemma-derived-quasi-coherent} の (2)(c) により，対象
$Lg_!K$ は $D_{\textit{LQCoh}^{fbc}}(\mathcal{O}_\mathcal{X})$ に属する．
ここから，$\mathcal{O}(x') \to \mathcal{O}(x)$ が平坦な環準同型なら
補題に表示した写像が同型であることは容易に従う．詳細は省略する．

\medskip\noindent
この段落では，問題が \'etale 位相に関して局所的であることを示す．
$x \to x'$ を $\mathcal{X}_{affine}$ の一般の射とする．
$\{x'_i \to x'\}$ を $\mathcal{X}_{affine, \etale}$ の被覆とする．
$x_i = x \times_{x'} x'_i$ と置けば，
$\{x_i \to x\}$ も $\mathcal{X}_{affine, \etale}$ の被覆である．
このとき $\mathcal{O}(x') \to \prod \mathcal{O}(x'_i)$ は
忠実平坦な \'etale 環準同型であり，
$$
\prod \mathcal{O}(x_i) =
\mathcal{O}(x) \otimes_{\mathcal{O}(x')} \left(\prod \mathcal{O}(x'_i)\right)
$$
である．したがって，省略する簡単な代数の議論により，補題の主張を
$\mathcal{X}_{affine}$ の各射 $x_i \to x'_i$ について証明すれば十分である．
言い換えれば，問題は \'etale 位相で局所的である．

\medskip\noindent
スキーム $X$ と全射的な滑らかな射 $f : X \to \mathcal{X}$ を選ぶ．
（記法の濫用により）$f$ を $\mathcal{X}$ の対象とみなしてよく，このとき
$(\Sch/X)_{fppf} = \mathcal{X}/f$ である．Sheaves on Stacks, Section
\ref{stacks-sheaves-section-restriction} を参照されたい．例えば
Sheaves on Stacks, Lemma
\ref{stacks-sheaves-lemma-surjective-flat-locally-finite-presentation} により，
$x'_i : U'_i = p(x'_i) \to \mathcal{X}$ が $f$ を経由するような
\'etale 被覆 $\{x'_i \to x'\}$ が存在する．前段落の結果により，
$x \to x'$ は，関手 $(\Sch/X)_{fppf} \to \mathcal{X}$ による
$(\textit{Aff}/X)_{fppf}$ の射 $U \to U'$ の像であると仮定してよい．
ここで $(\Sch/X)_{fppf}$ への $Lg_!K$ の制限は，Lemma
\ref{stacks-perfect-lemma-lisse-etale-functorial-derived} により
$f^*Lg_!K = L(g')_!(f')^*K$ に等しいことを用いる．
% P12-ERR-0050: Japanese removes the duplicated verb in the frozen phrase "we see use that".
これで次の段落で論じる場合に帰着される．

\medskip\noindent
$\mathcal{X} = (\Sch/X)_{fppf}$ であり，$x \to x'$ がアフィンスキームの射
$U \to U'$ に対応すると仮定する．なお $U'$ 上 \'etale（または Zariski）
局所的に作業できるので，$U' \to X$ が $X$ のあるアフィン開集合を
経由すると仮定してよい．これで次の段落で論じる場合に帰着される．

\medskip\noindent
$X = \Spec(R)$ がアフィンスキームであり，
$\mathcal{X} = (\Sch/X)_{fppf}$ で，$x \to x'$ がアフィンスキームの射
$U \to U'$ に対応すると仮定する．$R\Gamma(X, K)$ を表す
$R$-加群の複体を $M^\bullet$ とする．More on Algebra, Lemma
\ref{more-algebra-lemma-K-flat-resolution} の構成により，
各 $P_n^\bullet$ が自由 $R$-加群の上に有界な複体であるように
$M^\bullet = \colim P_n^\bullet$ と仮定してよい．詳細は省略する．
More on Algebra, Remark \ref{more-algebra-remark-P-resolution} も参照されたい．
$X_{flat, fppf} = (\Sch/X)_{flat, fppf}$ 上で，規則
$$
U \longmapsto \Gamma(U, M^\bullet \otimes_R \mathcal{O}_U)
$$
により与えられる加群の複体 $M^\bullet_{flat, fppf}$ を考える．
Descent, Section \ref{descent-section-quasi-coherent-sheaves} の議論により，
これは層の複体である．標準写像 $M^\bullet_{flat, fppf} \to K$ があり，
証明冒頭の所見により，これは $X_{flat, fppf}$ のアフィン対象上の
切断に同型を誘導する．$X_{flat, fppf}$ のすべての対象はアフィン対象に
よる被覆をもつので，$M^\bullet_{flat, fppf}$ は $K$ と一致する．

\medskip\noindent
$M^\bullet_{fppf}$ を，上と同じ表示式で与えられる $X_{fppf}$ 上の
加群の複体とする．$Lg_!\mathcal{O} = g_!\mathcal{O} = \mathcal{O}$ を
思い出そう．$Lg_!$ は $g_!$ の左導来関手なので，
$Lg_!P_{n, flat, fppf}^\bullet = P_{n, fppf}^\bullet$ である．
関手 $Lg_!$ はホモトピー余極限と可換し（または Cohomology on Sites,
Lemma \ref{sites-cohomology-lemma-existence-derived-lower-shriek} における
その構成により），$M^\bullet = \colim P_n^\bullet$ なので，
$Lg_!M^\bullet_{flat, fppf} = M^\bullet_{fppf}$ と結論される．
$U = \Spec(A)$，$U' = \Spec(A')$ と書き，$U \to U'$ が環準同型
$A' \to A$ に対応するとする．上から
$$
R\Gamma(U, Lg_!K) = M^\bullet \otimes_R A
\quad\text{and}\quad
R\Gamma(U', Lg_!K) = M^\bullet \otimes_R A'
$$
を得る．$M^\bullet$ は $R$-加群の K-flat 複体なので，テンソル積の
推移性により
$$
R\Gamma(U', Lg_!K) \otimes_{A'}^\mathbf{L} A
\longrightarrow
R\Gamma(U, Lg_!K)
$$
は求める擬同型である．
\end{proof}

\begin{proposition}
\label{stacks-perfect-proposition-QC-compare}
$\mathcal{X}$ を代数スタックとする．このとき
$\mathit{QC}(\mathcal{X})$ は $D_\QCoh(\mathcal{O}_\mathcal{X})$ と
標準的に同値である．
\end{proposition}

\begin{proof}
Sheaves on Stacks, Lemma \ref{stacks-sheaves-lemma-QC-compare-fppf} により，
比較射 $\epsilon : \mathcal{X}_{affine, fppf} \to \mathcal{X}_{affine}$
による引き戻しは，$\mathit{QC}(\mathcal{X})$ を充満部分圏
$Q_\mathcal{X} \subset D(\mathcal{X}_{affine, fppf}, \mathcal{O})$ と同一視する．
Sheaves on Stacks, Equation
(\ref{stacks-sheaves-equation-alternative-ringed}) における環付きトポスの
同値を用いて，$Q_\mathcal{X}$ を
$D(\mathcal{X}_{fppf}, \mathcal{O})$ の充満部分圏とみなし，以下でもそうする．

\medskip\noindent
同様に，Lemma \ref{stacks-perfect-lemma-bousfield-colocalization} および Remark
\ref{stacks-perfect-remark-QCoh-admissible} により，
$D_\QCoh(\mathcal{O}_\mathcal{X})$ は
$D_{\textit{LQCoh}^{fbc}}(\mathcal{O}_\mathcal{X})$ の左許容部分圏
$D_{\textit{Parasitic} \cap \textit{LQCoh}^{fbc}}(\mathcal{O}_\mathcal{X})$
の左直交補 $\mathcal{A}$ とみなすことができる．

\medskip\noindent
証明を完了するため，$D(\mathcal{X}_{fppf}, \mathcal{O})$ の部分圏として
$Q_\mathcal{X}$ が $\mathcal{A}$ に等しいことを示す．

\medskip\noindent
ステップ 1：$Q_\mathcal{X}$ は
$D_{\textit{LQCoh}^{fbc}}(\mathcal{O}_\mathcal{X})$ に含まれる．
$Q_\mathcal{X}$ の対象 $K$ は，
$D(\mathcal{X}_{affine, fppf}, \mathcal{O})$ の対象とみなした $K$ について
$R\epsilon_*K$ が $\mathit{QC}(\mathcal{X}_{affine}, \mathcal{O})$ の
対象であるという性質によって特徴づけられる．これはさらに，
$\mathcal{X}_{affine}$ のすべての射 $x \to x'$ に対して写像
$$
R\Gamma(x', K) \otimes_{\mathcal{O}(x')}^\mathbf{L} \mathcal{O}(x)
\longrightarrow
R\Gamma(x, K)
$$
が同型であることをちょうど意味する．Cohomology on Sites, Lemma
\ref{sites-cohomology-lemma-cartesion-plus-topology} の主張にある脚注を
参照されたい．いま，$x' \to x$ がアフィンスキームの平坦射上にあるなら，
% P12-ERR-0051: frozen source reverses x -> x' to x' -> x at the flat-map specialization; preserve formula literally.
これは
$$
H^i(x', K) \otimes_{\mathcal{O}(x')} \mathcal{O}(x)
\cong
H^i(x, K)
$$
を意味する．これは明らかに，$H^i(K)$ が \'etale 位相に対する層であり
（Sheaves on Stacks, Lemma
\ref{stacks-sheaves-lemma-quasi-coherent-alternative}），かつ平坦基底変換性を
もつことを意味する（細部は省略する）．

\medskip\noindent
ステップ 2：$Q_\mathcal{X}$ は $\mathcal{A}$ に含まれる．
これを示すには，$Q_\mathcal{X}$ の $K$ と
$D_{\textit{Parasitic} \cap \textit{LQCoh}^{fbc}}(\mathcal{O}_\mathcal{X})$
のすべての $P$ に対して $\Hom(K, P) = 0$ であることを示せば十分である．
対象
$$
H = R\SheafHom_{\mathcal{O}_\mathcal{X}}(K, P)
$$
を考える．$x$ を，アフィンスキーム $U = p(x)$ 上にある
$\mathcal{X}$ の対象とする．Cohomology on Sites, Lemma
\ref{sites-cohomology-lemma-section-RHom-over-U} により，次の第一の等式を得る：
$$
R\Gamma(x, H) =
R\Hom_{\mathcal{O}_\mathcal{X}}(K|_{\mathcal{X}/x}, P|_{\mathcal{X}/x}) =
R\Hom_{\mathcal{O}}(K|_{\mathcal{X}_{affine}/x}, P|_{\mathcal{X}_{affine}/x})
$$
第二の等式は，サイト $\mathcal{X}/x$ のトポスがサイト
$\mathcal{X}_{affine}/x$ のトポスと同値であることから従う．
Sheaves on Stacks, Equation
(\ref{stacks-sheaves-equation-alternative-ringed}) を参照されたい．
$K = \epsilon^*N$ と書ける．ここである $N$ は
$\mathit{QC}(\mathcal{O})$ に属する．すると Cohomology on Sites, Lemma
\ref{sites-cohomology-lemma-QC-hom-out-of-plus-topology} により
$$
R\Gamma(x, H) =
R\Hom_{D(\mathcal{O}(x))}(R\Gamma(x, N), R\Gamma(x, P))
$$
を得る．Lemma \ref{stacks-perfect-lemma-cohomology-parasitic} により，
$U \to \mathcal{X}$ が平坦なら $R\Gamma(x, P) = 0$ であり，したがって
同じ仮定のもとで $R\Gamma(x, H) = 0$ である．Lemma
\ref{stacks-perfect-lemma-cohomology-parasitic-complex} により
$R\Gamma(\mathcal{X}, H) = 0$ と結論され，ゆえに
$\Hom(K, P) = 0$ である．

\medskip\noindent
ステップ 3：$\mathcal{A}$ は $Q_\mathcal{X}$ に含まれる．
$K$ を $\mathcal{A}$ の対象，$x \to x'$ を
$\mathcal{X}_{affine}$ の射とする．写像
$$
R\Gamma(x', K) \otimes_{\mathcal{O}(x')}^\mathbf{L} \mathcal{O}(x)
\longrightarrow
R\Gamma(x, K)
$$
が擬同型であることを示さなければならない．Cohomology on Sites, Lemma
\ref{sites-cohomology-lemma-cartesion-plus-topology} の主張にある脚注を
参照されたい．Lemma \ref{stacks-perfect-lemma-bousfield-colocalization} の証明と Remark
\ref{stacks-perfect-remark-QCoh-admissible} の議論により，$\mathcal{A}$ は
$D_\QCoh(\mathcal{O}_{\mathcal{X}_{flat, fppf}})$ への $Lg_!$ の制限の像である．
したがって $K = Lg_!M$ と仮定してよい．ここである $M$ は
$D_\QCoh(\mathcal{O}_{\mathcal{X}_{flat, fppf}})$ に属する．
すると求める等式は Lemma
\ref{stacks-perfect-lemma-higher-shriek-QC} から従う．
% P12-ERR-0052: Japanese supplies singular agreement for the frozen "equality follow" wording.
\end{proof}
